\chapter{Problème statistique et chaîne de traitement}
\label{ch:problem}

\begin{lecture}
Le fichier fournit des élèves dont la maison est connue. Le modèle cherche des
régularités reliant leurs notes à leur maison, puis les applique à un élève
absent du fichier. Toute l'analyse porte sur l'écart entre la réussite sur les
élèves déjà vus et la réussite sur de nouveaux élèves.
\end{lecture}

On observe $m$ élèves. L'élève d'indice $i$ est décrit par $n$ variables numériques
$\x_i\in\R^n$ et par une maison $c_i$ appartenant à
\[
  \mathcal C=\{\text{Gryffindor},\text{Hufflepuff},
  \text{Ravenclaw},\text{Slytherin}\},\qquad K=|\mathcal C|=4.
\]
Le but est une règle $f:\R^n\to\mathcal C$ dont l'erreur reste faible sur des
élèves nouveaux, issus du même mécanisme de génération des données. La
performance sur l'échantillon observé sert d'indicateur intermédiaire.

Le terme \emph{régression} désigne ici la modélisation d'un log-rapport de chances
par une fonction affine. La variable réponse reste qualitative, ce qui range le
problème dans la classification \cite[chap.~4]{hastie2009}.

La table~\ref{tab:sample} montre quatre lignes du fichier. Ces valeurs passent par
l'imputation et la standardisation du chapitre~\ref{ch:data} avant d'entrer dans
le modèle. La colonne House porte la réponse à prédire.

\begin{table}[htbp]
\centering
\begin{tabular}{lrrrrl}
\toprule
élève & Astronomy & Herbology & Charms & Flying & House \\
\midrule
$1$ & $-487{,}9$ & $ 5{,}7$ & $-232{,}8$ & $-26{,}9$  & Ravenclaw \\
$2$ & $-552{,}1$ & $-6{,}0$ & $-252{,}2$ & $-113{,}5$ & Slytherin \\
$3$ & $-366{,}1$ & $ 7{,}7$ & $-227{,}3$ & $ 30{,}4$  & Ravenclaw \\
$4$ & $ 697{,}7$ & $-6{,}5$ & $-256{,}8$ & $200{,}6$  & Gryffindor \\
\bottomrule
\end{tabular}
\caption{Quatre lignes du jeu d'entraînement, réduites à quatre colonnes de notes
sur les treize disponibles. Une ligne représente une observation.}
\label{tab:sample}
\end{table}

\begin{figure}[htbp]
  \centering
  \resizebox{\textwidth}{!}{%
  \begin{tikzpicture}[
    node distance=5mm and 5mm,
    block/.style={draw=accent,rounded corners=1mm,fill=accentlight,
      align=center,minimum height=10mm,text width=27mm},
    arr/.style={-{Latex[length=2mm]},thick,draw=ink}]
    \node[block] (raw) {données\\brutes};
    \node[block,right=of raw] (split) {partition\\train--test};
    \node[block,right=of split] (prep) {imputation et\\standardisation};
    \node[block,right=of prep] (fit) {$K$ modèles\\binaires};
    \node[block,right=of fit] (eval) {décision et\\évaluation};
    \draw[arr] (raw)--(split); \draw[arr] (split)--(prep);
    \draw[arr] (prep)--(fit); \draw[arr] (fit)--(eval);
  \end{tikzpicture}}
  \caption{Chaîne de traitement. Les statistiques de préparation sont estimées sur
  l'ensemble d'apprentissage, puis appliquées telles quelles aux données de test.}
\label{fig:pipeline}
\end{figure}

\begin{workflow}{Ordre global : apprentissage, sélection et évaluation}
\begin{enumerate}
  \item \textbf{Partitionner.} Construire les indices d'apprentissage et de test,
        puis mettre le test à l'écart.
  \item \textbf{Préparer.} Estimer imputation, moyennes et écarts-types sur
        l'apprentissage ; transformer les données avec ces valeurs.
  \item \textbf{Organiser la sélection.} Réserver une validation ou construire
        des plis à l'intérieur des données disponibles pour l'apprentissage.
  \item \textbf{Former les cibles OvR.} Pour chaque classe $k$, construire le
        vecteur binaire $\y_k$.
  \item \textbf{Ajuster une configuration.} Pour chaque classe, initialiser
        $\thv_k$, puis répéter scores, probabilités, gradient et mise à jour
        jusqu'au critère d'arrêt.
  \item \textbf{Comparer les configurations.} Choisir $\alpha$, $\lambda$, le
        seuil et les autres hyperparamètres sur la seule validation.
  \item \textbf{Figer le modèle.} Conserver la préparation, l'ordre des
        caractéristiques, les $K$ vecteurs $\thv_k$ et l'ordre des classes.
  \item \textbf{Prédire le test.} Appliquer la préparation conservée, calculer
        les $K$ scores de chaque ligne et retenir leur maximum.
  \item \textbf{Évaluer une fois.} Comparer prédictions et classes réelles du
        test ; produire matrice de confusion et mesures par classe.
  \item \textbf{Archiver l'expérience.} Conserver graine, indices,
        hyperparamètres, historiques de convergence et résultats finaux.
\end{enumerate}
\end{workflow}

Chaque cadre « sous-procédure » du document développe une opération appelée par
cet ordre global : un calcul interne, un contrôle ou une méthode de sélection.

\begin{resultat}[Risque de généralisation]
La quantité pertinente est le \rsym{risque}{risque de généralisation}
$R(f)=\Pp(f(X)\neq C)$. La loi jointe de $(X,C)$ étant inconnue, on l'estime sur
un ensemble de test tenu à l'écart du choix des variables, de l'estimation des
transformations et de l'ajustement des paramètres.
\end{resultat}

Sur un test indépendant de $200$ élèves dont $160$ sont correctement classés,
l'erreur observée vaut $40/200=0{,}20$. Elle estime le risque de généralisation :
environ $20\,\%$ des futurs élèves comparables seraient mal classés. L'estimation
porte sa propre variabilité, un autre échantillon de $200$ élèves donnant un
résultat voisin et différent.

\section{Décision et calibration probabiliste}

Attribuer une maison à un nouvel élève est une décision. Estimer une probabilité
d'appartenance \rsym{calibration}{calibrée} est une tâche plus exigeante : parmi les élèves auxquels
le modèle attribue une probabilité proche de $0{,}8$, environ $80\,\%$
appartiennent effectivement à la classe. OvR répond à la première question ; la
seconde demande les précautions du chapitre~\ref{ch:ovr}.

Supposons que le modèle attribue une probabilité comprise entre $0{,}75$ et
$0{,}85$ à Ravenclaw pour $50$ élèves. Si $39$ sont réellement Ravenclaw, la
fréquence observée vaut $39/50=0{,}78$ et confirme l'annonce. Si $20$ le sont,
elle tombe à $0{,}40$ : le classement par maximum des scores reste utilisable, et
la lecture probabiliste de ces nombres demande une recalibration préalable.

\begin{resultat}[Objets définis par le problème]
Le jeu étiqueté fournit une matrice de caractéristiques et un vecteur de maisons.
L'apprentissage produit des paramètres. La prédiction transforme une ligne de
caractéristiques en une maison. L'évaluation compare le vecteur des maisons
prédites au vecteur des maisons réelles. Ces objets gardent des rôles séparés
dans toute la chaîne.
\end{resultat}
